%% GENERATED by python/paper/build.py from results/ -- do not edit by hand.
%% Rebuild:  .venv\Scripts\python.exe python\paper\build.py --only OECD_Countries
%% Source:   data/oecdFigure1.csv
\begin{table}[!htb]
\centering
\begin{threeparttable}
\caption{The cross-section of figure \ref{fig:US:OECDdata}, by country}
\label{table:OA:OECD_Countries}
\renewcommand{\arraystretch}{1.25}
\begin{tabularx}{\textwidth}{lYYYYYY}
\toprule
\textbf{Country} & \textbf{Pension spending, \% of GDP} & \textbf{Year} & \textbf{Population 2020 / 1990} & \textbf{Bismarckian index} & \textbf{Gini, disposable income} & \textbf{Year} \\
\midrule 
Austria & 14.3 & 2020 & 1.161 & 0.791 & 0.304 & 2020 \\ % row: aut
Belgium & 11.5 & 2020 & 1.155 & 0.616 & 0.274 & 2020 \\ % row: bel
Canada & 6.4 & 2020 & 1.373 & 0.595 & 0.318 & 2020 \\ % row: can
Czech Republic & 8.6 & 2020 & 1.033 & 0.603 & 0.254 & 2020 \\ % row: cze
Finland & 12.6 & 2020 & 1.109 & 0.944 & 0.277 & 2020 \\ % row: fin
\textbf{France} & 14.4 & 2020 & 1.192 & 0.849 & 0.321 & 2020 \\ % row: fra
Germany & 11.2 & 2020 & 1.047 & 0.778 & 0.323 & 2020 \\ % row: deu
Greece & 17.9 & 2020 & 1.049 & 0.857 & 0.331 & 2020 \\ % row: grc
Hungary & 7.7 & 2020 & 0.932 & 0.997 & 0.303 & 2020 \\ % row: hun
Ireland & 3.3 & 2020 & 1.435 & 0.500 & 0.302 & 2020 \\ % row: irl
Italy & 17.4 & 2020 & 1.048 & 0.956 & 0.358 & 2020 \\ % row: ita
Japan & 9.6 & 2020 & 1.020 & 0.747 & 0.330 & 2020 \\ % row: jpn
Luxembourg & 9.0 & 2020 & 1.651 & 0.848 & 0.341 & 2020 \\ % row: lux
New Zealand & 5.1 & 2020 & 1.496 & 0.501 & 0.320 & 2020 \\ % row: nzl
Norway & 7.6 & 2020 & 1.268 & 0.593 & 0.277 & 2020 \\ % row: nor
Poland & 11.4 & 2020 & 1.004 & 0.964 & 0.308 & 2020 \\ % row: pol
Portugal & 13.6 & 2020 & 1.040 & 0.892 & 0.328 & 2020 \\ % row: prt
Spain & 12.8 & 2020 & 1.218 & 0.877 & 0.343 & 2020 \\ % row: esp
Sweden & 8.5 & 2020 & 1.210 & 0.619 & 0.307 & 2020 \\ % row: swe
United Kingdom & 7.8 & 2020 & 1.166 & 0.550 & 0.329 & 2020 \\ % row: gbr
\textbf{United States} & 7.4 & 2020 & 1.328 & 0.687 & 0.419 & 2020 \\ % row: usa
\bottomrule
\end{tabularx}
\begin{tablenotes}[flushleft]
\footnotesize
\item[] \textit{Note:} The 21 OECD members plotted in figure \ref{fig:US:OECDdata}; the Bismarckian index is $1/(2-\theta)$ of the measure of the design in its note. Each variable is read in 2020 or the nearest year within 2 years of it (the year columns); population growth is the population of 2020 over that of 1990. Sources in table \ref{table:OA:OECD_Sources}. In bold the economies the paper calibrates.
\end{tablenotes}
\end{threeparttable}
\end{table}
